\section{Background}
\label{sec:background}

\subsection{Matrices, indices, and the public/internal split}

In \Lean{}, a matrix is the dependent function type \code{Matrix m n R}. The row
and column index types are parameters, so the public interface can be stated over
arbitrary finite types. This generality is convenient for users, but structural
recursion over an arbitrary finite type is awkward. The repository therefore
splits the development into two layers.

The kernel works over matrices whose row and column indices are both
\code{Fin}, and proves the recursive predicate
\code{IsHermiteNormalFormFin}. For arbitrary finite index types,
the structures \code{FiniteIndexing} and \code{MatrixIndexing} store explicit
equivalences to \code{Fin}. The public predicates are
\begin{center}
\code{IsHermiteNormalFormWith}\qquad
\code{IsHermiteNormalFormOrdered}.
\end{center}
They therefore state exactly which enumeration defines the normal-form
semantics. An \code{HNFResult} stores that indexing and the underlying strong
\code{Fin} result; matrices and certificates on the original index types are
derived projections.

\subsection{Euclidean computation versus canonical representatives}

The executable HNF layer assumes \code{EuclideanDomain R},
\code{NormalizationMonoid R}, and decidable equality. The Euclidean structure
supplies division, remainders, \(\gcd\), and the Bezout coefficients
\code{gcdA}/\code{gcdB} used in the elimination gadget. The normalization
structure supplies \code{normalize} and \code{normUnit}, which are enough to
state that pivots are chosen in a normalized associate class.

What \code{NormalizationMonoid} does \emph{not} provide is exact control over
the remainder operator \texttt{\%}. The computational interface of
\code{EuclideanDomain} lets the algorithm reduce entries modulo a pivot, but it
does not by itself guarantee two facts needed by the mathematical uniqueness
argument:

\begin{itemize}
  \item reducing an already reduced element is idempotent;
  \item two already reduced representatives modulo the same pivot are equal once
    their difference is divisible by that pivot.
\end{itemize}

The repository isolates exactly these requirements in the following
Prop-valued mixin.

\begin{lstlisting}[style=leanstyle,caption={The \code{CanonicalMod} boundary.},label={lst:canonicalmod}]
class CanonicalMod (R : Type _) [EuclideanDomain R] : Prop where
  mod_mod : forall a b : R, (a % b) % b = a % b
  reduced_eq_of_dvd_sub :
    forall {a b m : R}, a = a % m -> b = b % m -> m | b - a -> a = b
\end{lstlisting}

The field \code{mod_mod} closes the single-step reduction proof: when the top
row is explicitly replaced by a remainder, a second reduction must leave it
unchanged. The field \code{reduced_eq_of_dvd_sub} is stronger and is used
only in uniqueness: it upgrades congruence modulo a pivot into literal equality
once both sides are already canonical representatives. The current repository
provides instances for \code{Int} and \code{Polynomial Rat}. The development
remains parametric in \code{[CanonicalMod R]}.

\subsection{The recursive HNF predicate}

The internal specification \code{IsHermiteNormalFormFin} is the real HNF
definition. It is inductive, with cases for empty rows, empty columns, an all-zero
first column, and a pivot in the \((0,0)\) position. The pivot case packages the
usual row-style HNF conditions in a recursive form:

\begin{itemize}
  \item the top-left entry is nonzero;
  \item the top-left entry is normalized;
  \item every entry below the pivot in the first column is zero;
  \item the lower-right submatrix \(\lowerRight A\) is recursively in HNF;
  \item if the first nonzero entry of a lower row occurs in column \(j+1\), then
    the corresponding top-row entry is already the canonical remainder modulo
    that lower-row pivot.
\end{itemize}

The last clause is recorded by the field \code{hreduced} in the constructor
\code{pivot}. It is phrased through the computable search function
\code{firstNonzero?}, which returns the first nonzero position of a row or
\code{none}. This is the bridge between the mathematical statement ``entries
above later pivots are reduced'' and the executable reduction loop that modifies
only the top row.

This recursive specification is the row-oriented transpose of the HNF
conditions and Euclidean column-reduction schedule in Cohen's Algorithm 2.4.5
\citep[pp. 68--69]{cohen1993}, but in Lean it has two
decisive advantages. First, it matches the actual structure of the algorithm:
first isolate a pivot column, then recurse on a smaller matrix. Second, its
destructors expose exactly the local zero and reduction information that the
uniqueness proof needs.
